\section{Conclusions and open problems}\label{sec:conclusion}

Small treewidth makes a sparse mixed-integer problem tractable only together
with quantitative information about the objective. Two such quantities
suffice for certified global optimization on a box: upper bounds on the
coordinate curvature, which turn finite dynamic programming into a valid lower
bound and its min-marginals into a safe filter, and a quadratic growth
constant, which keeps the filtered grids at a size independent of the
accuracy. The resulting algorithms are fixed-parameter tractable in bag size
and condition number, need no knowledge of the growth constant, and produce
certificates that are checked by recomputation in exact arithmetic. The lower
bounds of Section~\ref{sec:limits} show that this parameterization is close to
the right one: neither parameter can be dropped, the dependence on the
condition number must be polynomial, and its exponent must grow with the bag
size. Exact recourse reduces the curvature that has to be discretized or
removes a residual problem from discretization altogether.

The certificates also have uses outside complete algorithms. A corrected-grid
bound on a sparse subproblem is a valid node bound inside spatial
branch-and-bound; the min-marginal filter is bound tightening whose
reductions are certified; and a replayable certificate allows the result of a
floating-point solver to be checked, which matters because such solvers satisfy
constraints only within tolerances (Section~\ref{sec:computation}). The
experiments show that the predicted grid sizes appear in practice and that the
width of the available decomposition, not the accuracy, is the binding limit.

Several questions remain open.
\begin{enumerate}[label=(\arabic*)]
\item \emph{Negative curvature.} Is box QP fixed-parameter tractable in the
width and $\nu/g$, where $\nu=\max\{0,-\lambda_{\min}(\nabla^2F)\}$?
Theorem~\ref{thm:cv} answers this when the convex part can be eliminated with
certified responses whose pieces are small, Proposition~\ref{prop:cv-limit}
shows that recourse alone cannot do it in general, and
Proposition~\ref{lim:prop:moments} shows that matching local moments of any
finite order is not enough. A positive answer would need curvature measured
only where near-optimal points can lie, or inequalities that couple both sides
of a separator.
\item \emph{Constants in the exponent.} The algorithm has
$f(p,\kappa)=(Cp\sqrt\kappa)^p\kappa^{O(1)}$, and the exponent of $\kappa$ must grow
linearly with $p$. Is the factor $p^p$, which comes from the $\log n$ growth of
the graded grids, necessary, and is the exponent $p/2$ of $\kappa$ optimal for
explicit inputs as it is for value oracles?
\item \emph{Coupling constraints.} Is there a feasible rounding scheme for
totally unimodular, or more general, sparse constraints that works with graded
grids? Propositions~\ref{prop:tu-misaligned} and the example after it show that
curvature-only corrections on graded grids fail, so a different correction is
needed.
\item \emph{Optimal sets.} When the optimal set is finite, can exact output be
obtained in time $f(p,\kappa_S,r)\poly(I)$, where $r$ bounds the optimal values per
coordinate? For optimal sets with continuous components, which compact
descriptions beyond the classes of Section~\ref{sec:optsets} can be computed
with a decomposition?
\item \emph{Exact polynomial output.} For continuous polynomial objectives the
minimizer can be irrational. Certificates that keep correlations between
coordinates, rather than signs over coordinate rectangles, seem necessary
under weak complementarity (Proposition~\ref{prop:weakcompl}); no construction of
parameterized size is known.
\end{enumerate}
